\documentclass[a4paper]{article}

\usepackage[intlimits]{amsmath}
\usepackage{amssymb}
\usepackage{pst-all}
\usepackage{pstricks-add}
\usepackage{multido}
\usepackage{calc}


\setlength{\hoffset}{2cm}
\addtolength{\hoffset}{-1in}
\addtolength{\hoffset}{-\oddsidemargin}
\setlength{\textwidth}{170mm}
\setlength{\voffset}{2cm}
\addtolength{\voffset}{-1in}
\addtolength{\voffset}{-\topmargin}
\addtolength{\voffset}{-\headheight}
\addtolength{\voffset}{-\headsep}
\setlength{\textheight}{257mm}

\pagestyle{empty}

\newcommand*{\hop}{\bigskip\noindent}

\begin{document}

\psset{unit=2.4cm}
\psset{linewidth=0.5pt}
\begin{center}
{\large{\bf A very nice function} (nose function)}

\hop
\(f(x)=\sqrt[3]{x^3+1}-x\).\qquad Analyse it!

\hop
Rotate it by 90 degrees to see a nose. :-)\\
\begin{pspicture}(-3.5,0)(3.5,3)
\psset{plotpoints=300}
\psplot{0}{3.5}{x 3 exp 1 add 1 3 div exp x sub}
\psplot{-3.5}{-1}{0 0 x sub 3 exp 1 sub 1 3 div exp sub x sub}
\psplot{-1}{0}{1 0 x sub 3 exp sub 1 3 div exp x sub}

%%%%%%%%%%%%%%%%%%%%%%%%%%%%%%
% add, div, mul, sub
% neg				  negates top number in stack
% mod				  remainder of division
% rand				  random integer in [0, 230]
% exp 				  Powers, not exponential
% dup				  This operate pushes a second copy of the topmost object on the operand stack.
% exch				  swaps top two elements in stack
% sin, cos, tan, acos, asin       in radians
% atan				  takes two arguments and returns the angle
% log, ln
% ceiling, floor, truncate, round
% sqrt                            square root
% abs                             absolute value
% fact                            for the factorial
% Sum                             for building sums
% IfTE                            for an easy case structure
%%%%%%%%%%%%%%%%%%%%%%%%%%%%%%

\psline{->}(-3.5,0)(3.5,0)
\uput{0.1}[315](3.5,0){\(x\)}
\uput{0.1}[45](0,2){\(\sqrt[3]{x^3+1}-x\)}
\psline{->}(0,-0.5)(0,2)
\psline(-0.1,1.587)(0.1,1.587)
\psline(-1,-0.1)(-1,0.1)
\uput{0.1}[0](0.1,1.587){\(2^{2/3}\)}
\uput{0.1}[270](-1,-0.1){\(-1\)}
\end{pspicture}
\end{center}
\end{document}
